\documentclass[11pt]{article}
\usepackage{wsi-style}
\wsirunhead{The Dissipative Decoder}
\title{Thermodynamic Cost Bounds on the Serial Decoding Throughput Basin}
\author{Grant Whitmer}
\date{March 2026}
\begin{document}
\maketitle

\begin{abstract}
Serial decoding systems---ribosome, cognition, language, and artificial intelligence---converge to 3--6 bits per processing event. This study models serial decoding as maximizing net information return---bits gained minus discrimination cost---and shows that the size of the optimal alphabet is set by how discrimination cost scales. Two cost regimes are identified. In Regime~A (alphabet-bound), biological pairwise molecular recognition makes discrimination expensive and quadratic in alphabet size; the net-return optimum $M^*=(\alpha\gamma\ln 2)^{-1/\alpha}$ then falls near $M = 20$ amino acids for the estimated biological cost parameters, placing biology inside the 3--6 bit basin. In Regime~B (capacity-bound), silicon discrimination is cheap and scales sub-linearly ($\alpha < 1$; a preliminary log-log fit across an energy benchmark gives a sub-linear exponent, with the full dataset to be released), so the same objective places the optimum at a far larger alphabet and silicon is not confined to the basin. The optimum is set by the discrimination-cost structure, not a parameter-free constant. The ribosome sits essentially on its informational (rate-distortion) floor---zero slack in bits---but thermodynamically it dissipates roughly $25\times$ the Landauer minimum ($\sim$80\,kT to encode $\sim$4.39 bits vs.\ a $\sim$3\,kT floor); a 7-billion-parameter transformer operates approximately $10^9$ above its Landauer floor. A temperature prediction tested across 29 organisms yields partial correlation $r = -0.451$ ($p = 0.014$) between amino acid entropy and growth temperature, controlling for GC-content. The basin is a speed limit for pairwise discriminators; silicon escapes it entirely. The two-regime framework resolves why AI lands in the biological throughput basin despite having no thermodynamic constraint on vocabulary size.
\end{abstract}

\noindent\textbf{Keywords:} dissipative decoding; discrimination energy; Landauer bound; kinetic proofreading; alphabet-bound regime; capacity-bound regime; ribosome efficiency; serial information processing; phase-space analysis; energy scaling

\section{Introduction}

The first four papers in this series established what the throughput basin is and where it sits. Thirty-one systems across six domains cluster between 3 and 6 bits per processing event~\cite{whitmer2026floor,whitmer2026basin}. The ribosome was predicted to 0.003 bits from pure physics~\cite{whitmer2026serial}. Evolutionary simulations independently reinvented the genetic code~\cite{whitmer2026basin}. But one question remained unanswered: why does the basin sit at 3--6 bits rather than 1--2 or 10--12?

This paper answers: energy. Every serial decoder faces a tradeoff between informational return (how much is learned per processing step) and discrimination cost (the energy required to reliably distinguish one symbol from another). Informational return grows logarithmically with alphabet size---diminishing returns. Discrimination cost depends on the physical mechanism of discrimination, and this is where biology and silicon fundamentally diverge.

\section{Materials and Methods}

\subsection{Discrimination Cost Function}

The general discrimination cost function has two components:
\begin{equation}
W(M, \varepsilon) = W_{\mathrm{alphabet}}(M) + W_{\mathrm{fidelity}}(\varepsilon)
\end{equation}
where $W_{\mathrm{alphabet}}(M)$ captures the cost of maintaining $M$ distinguishable symbols and $W_{\mathrm{fidelity}}(\varepsilon)$ captures the cost of achieving error rate $\varepsilon$. For pairwise molecular recognition (Regime~A), $W_{\mathrm{alphabet}}$ scales as $M(M-1)/2$---the number of pairwise distinctions required. For learned parameters (Regime~B), $W_{\mathrm{alphabet}}$ scales sub-linearly with model capacity.

\subsection{Two Cost Regimes}

\textbf{Regime~A (Alphabet-Bound, $\alpha > 1$):} Biological systems using pairwise molecular recognition. The ribosome distinguishes amino acids through dedicated aminoacyl-tRNA synthetases; each of the 21 amino acids has a dedicated molecular machine whose sole job is to recognize one amino acid and reject the other 20. The number of pairwise distinctions grows quadratically: $M(M-1)/2$. At $M = 21$, that is 210 pairwise comparisons. At $M = 30$, it would be 435---more than double the discrimination infrastructure for less than half a bit of additional information.

\textbf{Regime~B (Capacity-Bound, $\alpha < 1$):} Silicon AI systems using independent learned parameters. A transformer model performs matrix multiplication over independently learned parameters. The vocabulary is implemented as a cheap linear projection (softmax layer). Adding tokens to the vocabulary costs negligibly. The cost of adding parameters is sub-linear---each additional parameter costs less energy than the last.

\subsection{Net-Return Objective}

The quantity an efficient decoder maximizes is not raw efficiency (bits per joule) but the \emph{net} information retained after paying the discrimination cost:
\begin{equation}
U(M) = \log_2(M) - \gamma\, W(M) = \log_2(M) - \gamma M^\alpha,
\end{equation}
where $\gamma$ sets the energy scale of one unit of discrimination cost in bit-equivalents. Maximizing the bare ratio $\log_2(M)/(\gamma M^\alpha)$ is not the right objective: it peaks at $M^*=e^{1/\alpha}\approx 1.6$ for quadratic cost, degenerately favoring a two-symbol alphabet, and its peak location is independent of $\gamma$. For the net return, marginal information gain $1/(M\ln 2)$ falls with $M$ while marginal cost $\gamma\alpha M^{\alpha-1}$ rises for super-linear cost, and the optimum sits where they balance. Both regimes possess an optimum; what differs is \emph{where} it falls, and that is governed by the cost structure $(\alpha,\gamma)$---expensive, steeply scaling discrimination pulls the optimum to a small alphabet, cheap or sub-linearly scaling discrimination pushes it to a large one.

\subsection{Energy Benchmark: 27 AI Models}

Energy consumption was measured for 27 AI models on standardized hardware (NVIDIA RTX 5090) using \texttt{nvidia-smi} GPU power polling at 10\,ms intervals. Models spanned the Pythia family (14M to 1.4B parameters) and additional architectures. The relationship between energy and parameter count was assessed via log-log regression.

\subsection{Temperature Prediction: Kazusa Database}

To test whether discrimination cost increases with temperature (via the $kT$ term in the Landauer bound~\cite{landauer1961}), amino acid usage entropy was computed for 29 organisms spanning growth temperatures from 8 to 100$^\circ$C using the live Kazusa Codon Usage Database. GC-content was included as a control variable, as it is a known confound for temperature-driven compositional shifts~\cite{zeldovich2007}. The partial correlation between amino acid entropy and temperature, controlling for GC-content, was computed.

\section{Results}

\subsection{Phase-Space Analysis}

The net-return objective $U(M)=\log_2(M)-\gamma M^\alpha$ has a unique interior maximum for every $\alpha>0$, located at:
\begin{equation}
M^* = \left(\frac{1}{\alpha \gamma \ln 2}\right)^{1/\alpha}
\end{equation}
The optimum is therefore not parameter-free: it is fixed by the discrimination-cost structure $(\alpha,\gamma)$, and different substrates land in very different places. What separates the regimes is not the \emph{existence} of an optimum---both regimes have one---but \emph{where} it falls. For steep, expensive discrimination (large $\alpha$ and $\gamma$, as in pairwise molecular recognition), $M^*$ is small and lies inside the 3--6 bit basin. For cheap, sub-linearly scaling discrimination (small $\alpha$ and $\gamma$, as in learned silicon parameters), $M^*$ is large and lies far above the basin. There is no sharp $\alpha=1$ phase boundary; the optimum moves continuously as the cost structure changes. The empirically measurable question is therefore not ``super-linear or not?'' but ``how expensive, and how steeply scaling, is discrimination?''---a continuous property of the cost structure $(\alpha,\gamma)$ rather than a binary phase label.

\subsection{Silicon Energy Scaling}

Across the 27 models evaluated on the NVIDIA RTX 5090, silicon energy scales \emph{sub-linearly} with parameter count ($E \propto \mathrm{params}^{\alpha}$, $\alpha < 1$): each additional parameter costs less energy than the last. A preliminary log-log fit gives $\alpha \approx 0.9$; the precise exponent should be treated as provisional until the full 27-model dataset is released (the four disclosed Pythia points in Table~\ref{tab:pythia} alone fit a shallower slope), but the qualitative claim---sub-linearity---is what the argument requires and is robust. This is the opposite of biology's expensive quadratic scaling: cheap sub-linear cost places silicon's net-return optimum at a large alphabet, so AI is not confined to the 3--6 bit basin.

\begin{table}[ht]
\centering
\caption{Energy scaling across the Pythia model family (identical tokenizer).}
\begin{tabular}{lcccc}
\toprule
Model & Parameters & BPB & Energy/token (mJ) & Efficiency (bits/mJ) \\
\midrule
Pythia-14M & 14M & 1.74 & 0.12 & 14.5 \\
Pythia-70M & 70M & 1.23 & 0.18 & 6.8 \\
Pythia-410M & 410M & 0.97 & 0.34 & 2.9 \\
Pythia-1.4B & 1.4B & 0.85 & 0.58 & 1.5 \\
\bottomrule
\end{tabular}
\label{tab:pythia}
\end{table}

A 100-fold increase in parameters yields a 46\% improvement in compression but a nearly 5-fold increase in energy per token. Vocabulary size (identical across all Pythia models) contributes nothing to this trend---it is purely a capacity-driven scaling relationship, consistent with the Regime~B classification.

\subsection{Net-Return Optimum near $M \sim 20$}

For Regime~A with quadratic discrimination cost ($\alpha\approx 2$), the net-return optimum $M^*=(\alpha\gamma\ln 2)^{-1/\alpha}$ lands near $M=20$ when the cost coefficient takes a value of order $\gamma\approx1.8\times10^{-3}$. We do not derive $\gamma$ from first principles; it is calibrated so the optimum matches the observed alphabet, so this is a consistency result---the framework accommodates $M\approx20$ with a plausible cost coefficient---rather than a parameter-free prediction of it (see Limitations). On this reading the genetic code settled near 21 amino acids because that is where the marginal information from one more amino acid ($\sim$0.07 bit near $M=20$) no longer justifies its quadratic discrimination overhead---not because $M=20$ is a parameter-free constant. The 3--6 bit basin is the range these net-return optima occupy for expensive, super-linearly scaling discriminators.

\subsection{Landauer Comparison}

\begin{table}[ht]
\centering
\caption{Thermodynamic efficiency comparison.}
\begin{tabular}{lc}
\toprule
System & $\phi$ (ratio to Landauer minimum) \\
\midrule
Ribosome & ${\sim}25$ (${\sim}25\times$ above Landauer; ${\sim}1.02$ vs.\ the informational floor) \\
7B-parameter transformer & ${\sim}800{,}000{,}000$ ($10^9$ above minimum) \\
\bottomrule
\end{tabular}
\label{tab:landauer}
\end{table}

After 3.8 billion years of optimization, evolution has closed the \emph{informational} gap to essentially zero slack, while the ribosome's \emph{thermodynamic} dissipation remains roughly $25\times$ the Landauer minimum. Current AI systems operate approximately nine orders of magnitude above their Landauer floor.

\subsection{Slack Ordering}

The slack ordering across Regime~A systems is consistent with the thermodynamic framework: cheaper discrimination correlates with lower slack.

\begin{table}[ht]
\centering
\caption{Slack ordering across Regime~A systems.}
\begin{tabular}{lcccc}
\toprule
System & Floor $R_M(\varepsilon)$ & Observed & Slack $\Delta$ & Discrimination cost \\
\midrule
Ribosome & 4.39 & 4.39 & 0.00 & Low (GTP hydrolysis) \\
Phonology & 4.71 & 4.95 & 0.24 & Medium (neural) \\
Chromatic scale & 3.12 & 3.58 & 0.46 & Medium-high (perceptual) \\
Cognition & 1.97 & 3.12 & 1.15 & High (${\sim}10^4$ ATP/bit) \\
\bottomrule
\end{tabular}
\label{tab:slack}
\end{table}

Biology operates nearest the floor because molecular discrimination via GTP hydrolysis and tRNA selection is energetically cheap per comparison. Neural substrates operate with higher slack because synaptic discrimination costs approximately $10^4$ ATP per bit---orders of magnitude more expensive than molecular discrimination.

\subsection{Temperature Prediction}

Across 29 organisms spanning 8 to 100$^\circ$C, the partial correlation between amino acid entropy and growth temperature, controlling for GC-content, is $r = -0.451$ ($p = 0.014$). Hotter organisms use fewer distinct amino acids, consistent with the prediction that increased $kT$ raises the cost of each pairwise discrimination.

This result is preliminary, not proof. The GC-matched control comparisons were underpowered. Phylogenetic controls (PGLS) have not yet been applied. Prior work by Singer and Hickey (2003) and Zeldovich et al.\ (2007) documented temperature-driven compositional shifts attributed to GC bias and protein stability. The entropy measure captures a complementary signal, but the effect is moderate and requires confirmation.

\section{Discussion}

\subsection{Why Biology Stopped at 21 Amino Acids}

Biology could not decouple its vocabulary from its discrimination cost because molecular recognition is inherently pairwise: adding a new amino acid requires building entirely new molecular machinery to distinguish it from everything else. The cost grows quadratically while the information return grows logarithmically, so the net-return optimum for the estimated biological cost parameters sits near $M \sim 20$. This is why the genetic code stopped near 21 amino acids---it is the net-return optimum for these cost parameters, not a parameter-free constant.

\subsection{Why Silicon Escapes the Basin}

Silicon decoupled vocabulary from discrimination cost because learned parameters are independent. A softmax projection over a 50,000-token vocabulary is cheap regardless of token count. The cost scales sub-linearly with model capacity ($\alpha < 1$), so the net-return optimum sits at a large alphabet rather than near 20---silicon is not confined to the basin. AI can have 50,000 tokens for the same reason the ribosome cannot: cheap sub-linear discrimination and expensive quadratic discrimination place their optima in very different places.

\subsection{The Falsification That Became a Discovery}

We originally predicted that energy-efficient AI models would approach the rate-distortion floor, just as the ribosome does. The data showed the opposite: larger, more expensive models get closer to the floor. The prediction was cleanly falsified. But the falsification revealed something deeper than the original prediction would have: biology and silicon operate in fundamentally different cost regimes. The throughput basin is a feature of alphabet-bound systems specifically, not a universal law. This distinction---Regime~A versus Regime~B---is the paper's central contribution.

\subsection{Implications for AI Efficiency}

The nine-order-of-magnitude gap between AI and the Landauer floor suggests enormous room for improvement in silicon information processing efficiency. Current AI architectures waste the vast majority of their energy on overhead unrelated to the information-theoretic task. Whether this gap can be substantially closed through architectural innovations (neuromorphic computing, reversible computation) or whether it reflects fundamental constraints of digital hardware is an open question with significant practical implications.

\subsection{Dissipative Structure Interpretation}

Serial decoders in both regimes can be understood as dissipative structures in the sense of Prigogine: they maintain informational order by dissipating energy. The key difference is the scaling of dissipation with complexity. In Regime~A, dissipation scales super-linearly with alphabet size (the number of pairwise molecular comparisons), creating a thermodynamic ceiling that biology cannot exceed without fundamentally restructuring its discrimination mechanism. In Regime~B, dissipation scales sub-linearly with capacity (the number of learned parameters), allowing silicon systems to scale indefinitely---though with diminishing efficiency.

The within-family Pythia analysis provides direct evidence for Regime~B scaling: across the Pythia model series (14M to 1.4B parameters), BPB drops 46\% while using an identical tokenizer. A 100-fold increase in model parameters produces a nearly 5-fold increase in energy per token (Table~\ref{tab:pythia}, 0.12$\to$0.58 mJ) but only a 46\% improvement in compression. This is precisely the monotonically decreasing efficiency characteristic of Regime~B: larger models are better but less efficient.

\subsection{Implications for Synthetic Biology}

The net-return optimum near $M \sim 20$ has direct implications for efforts to expand the genetic code beyond the natural 20 amino acids. Adding a 22nd amino acid would require a new aminoacyl-tRNA synthetase capable of discriminating against all 21 existing amino acids---21 new pairwise comparisons---for approximately 0.07 additional bits of information per codon. The thermodynamic cost-benefit ratio becomes increasingly unfavorable as $M$ increases beyond 20. This analysis predicts that synthetic organisms with expanded genetic codes will exhibit reduced growth rates proportional to the discrimination cost overhead of the additional amino acids.

\subsection{Limitations}

\begin{enumerate}
\item The cost parameters ($\alpha$, $\gamma$) for biological substrates are estimates derived from biophysical literature, not direct measurements of discrimination cost per alphabet expansion.
\item The temperature prediction requires phylogenetic controls (PGLS) that have not yet been applied.
\item The energy benchmark covers only one GPU architecture; generalization to other hardware requires additional measurement.
\item The Landauer comparison assumes minimum-energy reversible computation as the baseline, which may not be achievable in practice.
\end{enumerate}

\section{Conclusion}

The throughput basin exists for biology because pairwise molecular recognition imposes expensive, quadratic discrimination cost, so the net-return optimum (logarithmic information return minus quadratic cost) falls in the 3--6 bit range for the estimated biological cost parameters. Silicon AI is not confined to this basin because matrix multiplication imposes cheap, sub-linear cost ($\alpha < 1$), which places its net-return optimum at a much larger alphabet. The distinction between Regime~A (alphabet-bound, expensive super-linear cost, optimum in the basin) and Regime~B (capacity-bound, cheap sub-linear cost, optimum far above it) resolves the apparent paradox of why AI converges to the biological throughput neighborhood despite having no thermodynamic constraint on vocabulary size. Preliminary temperature data from 29 organisms supports the prediction that increased thermal energy raises discrimination cost and reduces amino acid diversity. The basin is not a universal speed limit---it is a speed limit for pairwise discriminators. The companion paper~\cite{whitmer2026inherited} examines how this biological constraint propagates through human language and into AI training data, explaining the apparent paradox of why AI models converge to approximately 4.2 bits per token despite having no thermodynamic reason to do so. The two-regime framework developed here provides the thermodynamic foundation for that argument: Regime~A constrains biology directly through energy costs, and that constraint is then transmitted indirectly through the information structure of human language to AI systems operating in Regime~B. That AI inherits its throughput from the biological systems that generated its training data completes the causal chain from physics, through biology and cognition, through language, to artificial intelligence.

\subsection*{Author Contributions}

Conceptualization, G.L.W.; methodology, G.L.W.; software, G.L.W.; validation, G.L.W.; formal analysis, G.L.W.; investigation, G.L.W.; resources, G.L.W.; data curation, G.L.W.; writing---original draft preparation, G.L.W.; writing---review and editing, G.L.W.; visualization, G.L.W.; supervision, G.L.W.; project administration, G.L.W.

\subsection*{Funding}

This research received no external funding. All work was self-funded by the author.

\subsection*{Data Availability Statement}

The companion repository (\url{https://github.com/Windstorm-Labs/dissipative-decoder}) currently archives the manuscript only; the energy-benchmarking scripts, the 27-model dataset, and the temperature-analysis code are not yet deposited and will be released in a follow-up archive. The four Pythia data points in Table~\ref{tab:pythia} are reproducible from the table itself; the sub-linear exponent (preliminary value $\alpha\approx 0.9$) is fit to the full 27-model set, which will be included in the follow-up archive. Temperature data were obtained from the publicly available Kazusa Codon Usage Database.

\subsection*{Acknowledgments}

Mathematical derivations, energy benchmarking, and data analysis were performed with the assistance of Claude (Anthropic), an AI research tool. Temperature data were obtained from the publicly available Kazusa Codon Usage Database.

\subsection*{Conflicts of Interest}

The author declares no conflict of interest.

\subsection*{Data, Code \& Resources}
\noindent Manuscript source and archived record: \url{https://github.com/Windstorm-Institute/dissipative-decoder} (Zenodo: \href{https://doi.org/10.5281/zenodo.19432785}{10.5281/zenodo.19432785}). Experiments, datasets, and analysis code: \url{https://github.com/Windstorm-Labs/dissipative-decoder}. Windstorm Institute: \url{https://windstorminstitute.org}. Windstorm Labs: \url{https://windstormlabs.com}.

\begin{thebibliography}{99}

\bibitem{whitmer2026floor}
Whitmer, G.L., III.
The Receiver-Limited Floor: Rate-Distortion Bounds on Serial Decoding Throughput.
\emph{Zenodo}, 2026.
\href{https://doi.org/10.5281/zenodo.19322972}{DOI: 10.5281/zenodo.19322972}

\bibitem{whitmer2026basin}
Whitmer, G.L., III.
The Throughput Basin: Cross-Substrate Convergence and Decomposition of Serial Decoding Throughput.
\emph{Zenodo}, 2026.
\href{https://doi.org/10.5281/zenodo.19323193}{DOI: 10.5281/zenodo.19323193}

\bibitem{whitmer2026serial}
Whitmer, G.L., III.
The Serial Decoding Basin: Five Experiments on Convergence, Thermodynamic Anchoring, and Receiver-Limited Geometry.
\emph{Zenodo}, 2026.
\href{https://doi.org/10.5281/zenodo.19323422}{DOI: 10.5281/zenodo.19323422}

\bibitem{landauer1961}
Landauer, R.
Irreversibility and Heat Generation in the Computing Process.
\emph{IBM J.\ Res.\ Dev.}, 1961, \emph{5}, 183--191.
\href{https://doi.org/10.1147/rd.53.0183}{DOI: 10.1147/rd.53.0183}

\bibitem{zeldovich2007}
Zeldovich, K.B.; Berezovsky, I.N.; Shakhnovich, E.I.
Protein and DNA Sequence Determinants of Thermophilic Adaptation.
\emph{PLoS Comput.\ Biol.}, 2007, \emph{3}, e5.
\href{https://doi.org/10.1371/journal.pcbi.0030005}{DOI: 10.1371/journal.pcbi.0030005}

\bibitem{whitmer2026inherited}
Whitmer, G.L., III.
The Inherited Constraint: Biological Throughput Limits Shape the Information Structure of Human Language and AI.
\emph{Zenodo}, 2026.
\href{https://doi.org/10.5281/zenodo.19432910}{DOI: 10.5281/zenodo.19432910}

\bibitem{hopfield1974}
Hopfield, J.J.
Kinetic Proofreading: A New Mechanism for Reducing Errors in Biosynthetic Processes.
\emph{Proc.\ Natl.\ Acad.\ Sci.\ USA}, 1974, \emph{71}, 4135--4139.
\href{https://doi.org/10.1073/pnas.71.10.4135}{DOI: 10.1073/pnas.71.10.4135}

\bibitem{shannon1948}
Shannon, C.E.
A Mathematical Theory of Communication.
\emph{Bell Syst.\ Tech.\ J.}, 1948, \emph{27}, 379--423.
\href{https://doi.org/10.1002/j.1538-7305.1948.tb01338.x}{DOI: 10.1002/j.1538-7305.1948.tb01338.x}

\bibitem{whitmer2026fons}
Whitmer, G.L., III.
The Fons Constraint: Information-Theoretic Convergence on Encoding Depth in Self-Replicating Systems.
\emph{Zenodo}, 2026.
\href{https://doi.org/10.5281/zenodo.19274047}{DOI: 10.5281/zenodo.19274047}

\bibitem{whitmer2026code}
Whitmer, G.L., III.
Dissipative Decoder: Experiment Code and Data.
\emph{GitHub}, 2026.
Available online: \url{https://github.com/Windstorm-Labs/dissipative-decoder}

\end{thebibliography}

\bigskip
\noindent\emph{This is Paper~5 of The Windstorm Series. The series includes Paper~1: The Fons Constraint~\cite{whitmer2026fons}, Paper~2: The Receiver-Limited Floor~\cite{whitmer2026floor}, Paper~3: The Throughput Basin~\cite{whitmer2026basin}, Paper~4: The Serial Decoding Basin~\cite{whitmer2026serial}, Paper~6: The Inherited Constraint~\cite{whitmer2026inherited}, and Paper~7: The Throughput Basin Origin (DOI: 10.5281/zenodo.19498582).}

\bigskip
\noindent\textcopyright{} 2026 Grant Whitmer. This article is licensed under a Creative Commons Attribution 4.0 International License (CC BY 4.0).

\end{document}
